\subsection{Stretching the space}
Helmholtz equation
\begin{equation}\label{helmholtz}
	(\Delta +k^2)f=(\partial_x^2+\partial_y^2+k^2)f=0
\end{equation}

Consider the following change of variables
\begin{equation*}
	u=3x \qquad
	v=3y,
\end{equation*}
what form does \eqref{helmholtz} take in the \(u,v\) coordinate system?

Using the chain rule we get:
\begin{equation*}
	\partial_x=\dd{u}{x}\partial_u=3\partial_u \qquad
	\partial_y=\dd{v}{y}\partial_v=3\partial_v,
\end{equation*}
replacing this in \eqref{helmholtz} we get
\begin{equation*}
	(9\partial_u^2+9\partial_v^2+k^2)f=0,
\end{equation*}
dividing by 9 we obtain
\begin{equation*}
	(\partial_u^2+\partial_v^2+(k/3)^2)f=0,
\end{equation*}
hence it is equivalent to solving the original Helmholtz equation
\eqref{helmholtz} with \(k\) replaced by \(k/3\). In other words, stretching
the space by a factor of 3 is equivalent to dividing \(k\) by a factor of 3.
Here we see that changing the composition of the medium and changing the
geometry that light perceives is equivalent.

\subsection{Inverting with respect to the unit circle}
The coordinate transformation we will now use is

\begin{equation*}
	u=\frac{x}{x^2+y^2} \qquad
	v=\frac{y}{x^2+y^2}.
\end{equation*}
This is inversion with respect to the unit circle.

Following the same procedure as in the previous example we get
\begin{equation*}
	\partial_x=\dd{u}{x}\partial_u + \dd{v}{x}\partial_v \qquad
	\partial_y=\dd{u}{y}\partial_u + \dd{v}{y}\partial_v,
\end{equation*}
then
\begin{align*}
	\partial_x^2=
	  \left(\dd{u}{x}\right)^2\partial_u^2
	+ \left(\dd{v}{x}\right)^2\partial_v^2
	+ 2\dd{u}{x}\dd{v}{x}\partial_u\partial_v \\
	\partial_y^2=
	  \left(\dd{u}{y}\right)^2\partial_u^2
	+ \left(\dd{v}{y}\right)^2\partial_v^2
	+ 2\dd{u}{y}\dd{v}{y}\partial_u\partial_v
\end{align*}
and so
\begin{align*}
	\Delta &= \partial_x^2+\partial_y^2\\
	       &= \left(\dd{u}{x}\right)^2\partial_u^2
	+ \left(\dd{v}{x}\right)^2\partial_v^2
	+ 2\dd{u}{x}\dd{v}{x}\partial_u\partial_v\\
	       &+
	  \left(\dd{u}{y}\right)^2\partial_u^2
	+ \left(\dd{v}{y}\right)^2\partial_v^2
	+ 2\dd{u}{y}\dd{v}{y}\partial_u\partial_v\\
	       &=\left(\left(\dd{u}{x}\right)^2
	+\left(\dd{u}{y}\right)^2\right)\partial_u^2
	+\left(\left(\dd{v}{x}\right)^2
	+\left(\dd{v}{y}\right)^2\right)\partial_v^2\\
	       &+2\left(\dd{u}{x}\dd{v}{x}
	+\dd{u}{y}\dd{v}{y}\right)\partial_u\partial_v
\end{align*}
i.e.
\begin{multline}\label{laplacianuv}
	\Delta =\left(\left(\dd{u}{x}\right)^2
	+\left(\dd{u}{y}\right)^2\right)\partial_u^2
	       +\left(\left(\dd{v}{x}\right)^2
	+\left(\dd{v}{y}\right)^2\right)\partial_v^2\\
	       +2\left(\dd{u}{x}\dd{v}{x}
	+\dd{u}{y}\dd{v}{y}\right)\partial_u\partial_v
\end{multline}
In this case
\begin{equation}\label{partialu}
	\dd{u}{x}=\frac{y^2-x^2}{(x^2+y^2)^2} \qquad
	\dd{u}{y}=-\frac{2xy}{(x^2+y^2)^2},
\end{equation}
and
\begin{equation*}
	\dd{v}{x}=-\frac{2xy}{(x^2+y^2)^2} \qquad
	\dd{v}{y}=\frac{x^2-y^2}{(x^2+y^2)^2}.
\end{equation*}

Notice that
\begin{equation}\label{partialsRelations}
	\dd{v}{y}=-\dd{u}{x} \qquad
	\dd{v}{x}=\dd{u}{y}.
\end{equation}

Using \eqref{partialsRelations} to get rid of the \(v\) variable in \eqref{laplacianuv} we get
\begin{equation}\label{laplacianuvsimplified}
	\Delta=\left(
	\left(\dd{u}{x}\right)^2
	+\left(\dd{u}{y}\right)^2\right)
	\left(\partial_u^2+\partial_v^2 \right).
\end{equation}
Going back to the Helmholtz equation if we want to write it in the \(u,v\)
coordinates we get
\begin{equation*}
\left(
\left(\left(\dd{u}{x}\right)^2
+\left(\dd{u}{y}\right)^2\right)
\left(\partial_u^2+\partial_v^2\right)+ k^2\right)f=0
\end{equation*}
and dividing by \(\left(\dd{u}{x}\right)^2
+\left(\dd{u}{y}\right)^2\) we obtain:
\begin{equation}
\left(
	\partial_u^2+\partial_v^2
	+\left(
	k\middle/\sqrt{\left(\dd{u}{x}\right)^2
	+\left(\dd{u}{y}\right)^2}
	\right)^2
\right)f=0,
\end{equation}
so once again we can achieve the transformation just by changing the
parameter \(k\). In this example the parameter we need is a lot more
complicated than the one in the previous example; however, we can see that it
still is just a function of space, as it depends only on \(x\) and \(y\) as
seen in \eqref{partialu}.
\subsection{A non-example}
Not all transformations can be achieved by modifying \(k\), for example consider
\begin{equation*}
	u=2x \qquad
	v=3y,
\end{equation*}
In this case
\begin{equation*}
	\partial_x=2\partial_u \qquad
	\partial_y=3\partial_v
\end{equation*}
and then
\begin{equation*}
	(4\partial_u^2 + 9\partial_v^2 + k^2)f=0,
\end{equation*}
so we can not rewrite this equation to have the same form as the original
Helmholtz equation.
